\documentclass[8pt,a4paper]{extarticle}
\usepackage[a4paper,margin=0.42in,top=0.43in,bottom=0.43in]{geometry}
\usepackage[T1]{fontenc}
\usepackage{lmodern,multicol,listings,xcolor,enumitem}
\pagestyle{empty}
\setlength{\parindent}{0pt}
\setlength{\parskip}{1pt}
\setlength{\columnsep}{0.23in}
\setlist{nosep,leftmargin=*}
\lstset{language=C,basicstyle=\ttfamily\fontsize{6.3}{7.0}\selectfont,breaklines=true,breakatwhitespace=false,columns=fullflexible,keepspaces=true,showstringspaces=false,aboveskip=2pt,belowskip=2pt}
\newcommand{\heading}[1]{\begin{center}{\bfseries\large GRAPHIC ERA DEEMED TO BE UNIVERSITY}\\[-1pt]{\bfseries Department of Computer Science and Engineering}\\[-1pt]{\bfseries OS LAB -- MID-SEMESTER PRACTICAL: SET #1 SOLUTIONS}\end{center}\vspace{-5pt}\hrule\vspace{3pt}}
\newcommand{\prog}[1]{\vspace{3pt}{\bfseries #1}\par}
\newcommand{\answers}{\vspace{3pt}{\bfseries Short answers (4 $\times$ 2 marks)}\par}
\begin{document}
\heading{1}
\begin{multicols}{2}
\prog{P1 (6 marks): Parent odd sum; child even sum}
\begin{lstlisting}
#include <stdio.h>
#include <stdlib.h>
#include <unistd.h>
#include <sys/wait.h>
int main(void) {
 int a[100], n, i, sum=0;
 pid_t p;
 printf("Size (1..100): ");
 if(scanf("%d",&n)!=1||n<1||n>100)
   return 1;
 for(i=0;i<n;i++) scanf("%d",&a[i]);
 fflush(stdout);
 p=fork();
 if(p<0){perror("fork");return 1;}
 if(p==0){
   for(i=0;i<n;i++)
     if(a[i]%2==0) sum+=a[i];
   printf("Child PID=%ld PPID=%ld"
          " even sum=%d\n",
          (long)getpid(),
          (long)getppid(),sum);
   exit(0);
 }
 for(i=0;i<n;i++)
   if(a[i]%2!=0) sum+=a[i];
 printf("Parent PID=%ld PPID=%ld"
        " odd sum=%d\n",
        (long)getpid(),
        (long)getppid(),sum);
 waitpid(p,NULL,0);
 return 0;
}
\end{lstlisting}
\textbf{Example:} Input \texttt{1 2 3 4 5}; parent odd sum $=9$, child even sum $=6$. Output order may vary.
\prog{P2 (6 marks): Child sends string to parent through pipe}
\begin{lstlisting}
#include <stdio.h>
#include <stdlib.h>
#include <string.h>
#include <unistd.h>
#include <sys/wait.h>
int main(void) {
 int fd[2]; pid_t p; char s[100]={0};
 if(pipe(fd)==-1){perror("pipe");return 1;}
 p=fork();
 if(p<0){perror("fork");return 1;}
 if(p==0){
   close(fd[0]);
   const char *msg="Hello from child";
   if(write(fd[1],msg,strlen(msg)+1)<0)
     perror("write");
   close(fd[1]); exit(0);
 }
 close(fd[1]);
 ssize_t k=read(fd[0],s,sizeof(s)-1);
 if(k<0) perror("read");
 else {s[k]='\0';printf("Parent: %s\n",s);}
 close(fd[0]); waitpid(p,NULL,0);
 return k<0 ? 1:0;
}
\end{lstlisting}
\columnbreak
\answers
\textbf{Q1.} \texttt{fork()} returns $0$ in the child, the child's PID ($>0$) in the parent, and $-1$ on failure (no child created).\\
\textbf{Q2.} \texttt{wait()} blocks the parent until a child changes to a waitable state, normally termination; it retrieves the status and reaps the child, avoiding a zombie.\\
\textbf{Q3.} A \emph{zombie} has terminated but its parent has not reaped its status. An \emph{orphan} is still running after its parent exits and is reparented to an appropriate reaper.\\
\textbf{Q4.} \texttt{fd[0]} is the read end and \texttt{fd[1]} is the write end of an unnamed pipe.
\vspace{5pt}\hrule\vspace{3pt}
\textbf{Compile:} \texttt{gcc file.c -o app}\quad \textbf{Run:} \texttt{./app}.\\
\textbf{Marking per program (6):} correct setup 2, core logic 2, IPC/synchronization 1, output/cleanup 1.
\end{multicols}
\newpage
\heading{2}
\begin{multicols}{2}
\prog{P1 (6 marks): Child 1--10; parent 11--20, synchronized}
\begin{lstlisting}
#include <stdio.h>
#include <stdlib.h>
#include <unistd.h>
#include <sys/wait.h>
int main(void) {
 pid_t p=fork(); int i;
 if(p<0){perror("fork");return 1;}
 if(p==0){
   for(i=1;i<=10;i++) printf("%d ",i);
   printf("\n"); fflush(stdout);
   _exit(0);
 }
 if(waitpid(p,NULL,0)<0){
   perror("waitpid");return 1;
 }
 for(i=11;i<=20;i++) printf("%d ",i);
 printf("\n");
 return 0;
}
\end{lstlisting}
\textbf{Expected order:} Child prints \texttt{1 2 ... 10}, then parent prints \texttt{11 12 ... 20}. \texttt{waitpid()} enforces this order.
\prog{P2 (6 marks): Parent sends integer; child prints square}
\begin{lstlisting}
#include <stdio.h>
#include <stdlib.h>
#include <unistd.h>
#include <sys/wait.h>
int main(void) {
 int fd[2], n; pid_t p;
 printf("Integer: ");
 if(scanf("%d",&n)!=1) return 1;
 fflush(stdout);
 if(pipe(fd)<0){perror("pipe");return 1;}
 p=fork();
 if(p<0){perror("fork");return 1;}
 if(p==0){
   int x; close(fd[1]);
   ssize_t r=read(fd[0],&x,sizeof x);
   close(fd[0]);
   if(r!=sizeof x){
     fprintf(stderr,"read failed\n");
     _exit(1);
   }
   printf("Square = %lld\n",
          (long long)x*x);
   fflush(stdout); _exit(0);
 }
 close(fd[0]);
 if(write(fd[1],&n,sizeof n)!=sizeof n)
   perror("write");
 close(fd[1]); waitpid(p,NULL,0);
 return 0;
}
\end{lstlisting}
\textbf{Example:} Input \texttt{7}; child output \texttt{Square = 49}.
\columnbreak
\answers
\textbf{Q1.} \texttt{getpid()} returns the calling process's PID; \texttt{getppid()} returns its parent process's PID.\\
\textbf{Q2.} With ordinary \texttt{wait()}, the parent blocks until a child is waitable, typically after the child terminates; it then collects the child's status.\\
\textbf{Q3.} An unnamed pipe is a kernel-managed byte stream identified by file descriptors, not a pathname. Creating it before \texttt{fork()} allows the child to inherit descriptors for the same pipe.\\
\textbf{Q4.} If the pipe is empty but some write end is open, blocking \texttt{read()} waits for data or EOF. If all write ends close and no data remains, \texttt{read()} returns \texttt{0} (EOF).
\vspace{5pt}\hrule\vspace{3pt}
\textbf{Compile:} \texttt{gcc file.c -o app}\quad \textbf{Run:} \texttt{./app}.\\
\textbf{Marking per program (6):} process/pipe creation 2, computation 2, synchronization/descriptor handling 1, output 1.
\end{multicols}
\newpage
\heading{3 -- ADVANCED}
\begin{multicols}{2}
\prog{P1 (6 marks): Exactly three children, three \texttt{fork()} calls (no process-creation loop)}
\begin{lstlisting}
#include <stdio.h>
#include <stdlib.h>
#include <unistd.h>
#include <sys/wait.h>
int prime(int x){
 if(x<2)return 0;
 for(int d=2;d<=x/d;d++)
   if(x%d==0)return 0;
 return 1;
}
int square(int x){
 if(x<0)return 0;
 for(int d=0;(long long)d*d<=x;d++)
   if(d*d==x)return 1;
 return 0;
}
int main(void){
 int a[100],n,s=0,st;
 pid_t p1,p2,p3;
 printf("Size: ");
 if(scanf("%d",&n)!=1||n<1||n>100)
   return 1;
 for(int i=0;i<n;i++)scanf("%d",&a[i]);
 fflush(stdout);
 p1=fork();
 if(p1<0){perror("fork1");return 1;}
 if(p1==0){
   for(int i=0;i<n;i++)
     if(prime(a[i]))s+=a[i];
   printf("C1 PID=%ld PPID=%ld"
          " prime sum=%d\n",
          (long)getpid(),(long)getppid(),s);
   fflush(stdout);_exit(0);
 }
 p2=fork();
 if(p2<0){perror("fork2");
   waitpid(p1,NULL,0);return 1;}
 if(p2==0){
   for(int i=0;i<n;i++)
     if(square(a[i]))s+=a[i];
   printf("C2 PID=%ld PPID=%ld"
          " square sum=%d\n",
          (long)getpid(),(long)getppid(),s);
   fflush(stdout);_exit(0);
 }
 p3=fork();
 if(p3<0){perror("fork3");
   waitpid(p1,NULL,0);
   waitpid(p2,NULL,0);return 1;}
 if(p3==0){
   for(int i=0;i<n;i++)
     if(a[i]%3==0&&a[i]%5==0)s+=a[i];
   printf("C3 PID=%ld PPID=%ld"
          " divisible-by-15 sum=%d\n",
          (long)getpid(),(long)getppid(),s);
   fflush(stdout);_exit(0);
 }
 pid_t p[3]={p1,p2,p3};
 for(int i=0;i<3;i++){
   if(waitpid(p[i],&st,0)<0)perror("waitpid");
   else printf("Child %d exit=%d\n",i+1,
     WIFEXITED(st)?WEXITSTATUS(st):-1);
 }
 puts("Parent: all children finished");
 return 0;
}
\end{lstlisting}
\textbf{Note:} Each child exits before the next \texttt{fork()}, so only the original parent executes all three forks. A loop is used only for collecting statuses.
\columnbreak
\prog{P2 (6 marks): Two pipes, array request and result response}
\begin{lstlisting}
#include <stdio.h>
#include <stdlib.h>
#include <unistd.h>
#include <sys/wait.h>
typedef struct {int n,a[100];} Request;
typedef struct {int found,second;
                long long evenSum;} Reply;
int main(void){
 int toC[2],toP[2]; pid_t p;
 Request q={0}; Reply r={0};
 printf("Size (1..100): ");
 if(scanf("%d",&q.n)!=1||q.n<1||q.n>100)
   return 1;
 for(int i=0;i<q.n;i++)
   scanf("%d",&q.a[i]);
 fflush(stdout);
 if(pipe(toC)<0){perror("pipe");return 1;}
 if(pipe(toP)<0){perror("pipe");return 1;}
 p=fork();
 if(p<0){perror("fork");return 1;}
 if(p==0){
   close(toC[1]);close(toP[0]);
   if(read(toC[0],&q,sizeof q)!=sizeof q)
     _exit(1);
   int max=0,hasMax=0;
   for(int i=0;i<q.n;i++){
     int x=q.a[i];
     if(x%2==0)r.evenSum+=x;
     if(!hasMax||x>max){
       if(hasMax){r.second=max;r.found=1;}
       max=x;hasMax=1;
     }else if(x<max&&
              (!r.found||x>r.second)){
       r.second=x;r.found=1;
     }
   }
   if(write(toP[1],&r,sizeof r)!=sizeof r)
     _exit(1);
   close(toC[0]);close(toP[1]);
   _exit(0);
 }
 close(toC[0]);close(toP[1]);
 if(write(toC[1],&q,sizeof q)!=sizeof q)
   perror("write");
 close(toC[1]);
 if(read(toP[0],&r,sizeof r)==sizeof r){
   if(r.found)printf("Second distinct=%d\n",
                     r.second);
   else puts("No second distinct value");
   printf("Even sum=%lld\n",r.evenSum);
 }else perror("read");
 close(toP[0]);waitpid(p,NULL,0);
 return 0;
}
\end{lstlisting}
\answers
\textbf{Q1.} Three unconditional forks create $2^3=8$ total processes, of which $7$ are newly created children.\\
\textbf{Q2.} \texttt{waitpid(pid,\&status,WNOHANG)} returns $0$ if the specified child is still running and no status is available.\\
\textbf{Q3.} Empty pipe + open writer: blocking read waits. Empty pipe + all writers closed: read returns $0$ (EOF).\\
\textbf{Q4.} Retaining an unused write descriptor prevents EOF: a reader waiting for EOF may block indefinitely. Circular blocking can also occur when each process waits for input from the other.
\end{multicols}
\end{document}
